%%
%% TOMS Algorithm Article
%% Manuscript format for review
%%
\documentclass[manuscript,screen,review]{acmart}

%% Journal information for TOMS
\setcopyright{rightsretained}
\acmJournal{TOMS}
\acmYear{2026}
\acmVolume{1}
\acmNumber{1}
\acmArticle{1}
\acmMonth{1}
\acmDOI{10.1145/XXXXXXX}

%% BibTeX command
\AtBeginDocument{%
  \providecommand\BibTeX{{%
    Bib\TeX}}}

%% Bibliography
\bibliographystyle{ACM-Reference-Format}

\begin{document}

%%
%% Title
%%
\title{Algorithm XXX: Randompack, A Portable C Library for Reproducible Random Number Generation}

%%
%% Author information - FILL IN YOUR DETAILS
%%
\author{Kristján Jónasson}
\authornote{Corresponding author.}
\email{jonasson@hi.is}
\orcid{0000-0000-0000-0000}
\affiliation{%
  \institution{University of Iceland}
  \city{Reykjavík}
  \country{Iceland}
}


%%
%% Abstract
%%
\begin{abstract}
\end{abstract}

%%
%% Keywords
%%
\keywords{}

\maketitle

\section{Introduction}

The computer generation of pseudo-random numbers has a history dating back to the first
computers, with von Neumann's middle-square idea (1946)
\cite{vonneumann1951rejection,lecuyer2017history} and Lehmer's linear congruential
generator (LCG, 1949) \cite{lehmer1951mathematical,lecuyer2017history}. The first
generators were aimed at generating random integers (typically in $[0, 2^k)$) and
uniformly distributed reals in $[0,1]$ (through division by $2^k$), but soon other
distributions followed. In fact von Neumann's paper \cite{vonneumann1951rejection} had a
clever algorithm for drawing exponentially distributed samples, also obtainable by taking
the logarithm of uniform samples. An early and seminal method for random normals is given
by Box and Muller \cite{box1958normal}. Other important 20th century milestones are the
Tausworthe random number generator (RNG) \cite{tausworthe1965linear}, the Wichmann-Hill
RNG \cite{wichmann1982efficient}, the Park-Miller LCG RNG \cite{park1988random}, RANLUX
\cite{luscher1994portable} and the Mersenne Twister \cite{matsumoto1998mersenne}.

Parallel to the uniform generator development new methods for generating non-uniform
variates were developed. The normal distribution plays an important role, and other
important ones are the exponential and the gamma distributions. Random sampling from
several other distributions can be done by combining samples from these three (see Section
xxx). The already mentioned Box-Muller method (1958) requires the evaluation of two
logarithms, a sine, and a cosine for each pair of samples, but a few years later Marsaglia
and Bray proposed the polar method where the two trigonometrict functions were replaced by
a square root and rejection sampling with acceptance rate $\pi/4$. Rejection sampling had
already been used by von Neuman in his exponential method \cite{vonneumann1951rejection},
and it plays a central role in the ziggurat method for normal sampling, devised by
Marsaglia and Tsang in 


At the end of the 20th century and especially in the 21st century random number research
focus shifted to studying the statistical quality of RNGs, ensuring long periods, and
devising methods to create independent streams of random numbers, spurred on by faster,
often parallel computers. On a modern computer the Park-Miller RNG will start to repeat
itself after a few seconds on a single core.

Among the first studies of 

generation of
normally distributed

Among important milestones in this
history are the Park-Miller

\section{Design and interface}

\section{Algorithms and implementation}

\section{Evaluation}

\section{Discussion}

%%
%% Bibliography
%%
\bibliography{randompack}

%%
%% Appendix (optional)
%%
\appendix

\section{Algorithm Listings}

\section{Additional Performance Data}

\end{document}
